\documentclass[11pt,a4paper]{article}
\usepackage[margin=1.05in]{geometry}
\usepackage{amsmath,amssymb,amsthm}
\usepackage{booktabs}
\usepackage{microtype}
\usepackage[hidelinks]{hyperref}

\newtheorem{lemma}{Lemma}
\newtheorem{corollary}{Corollary}
\newtheorem{remark}{Remark}

\title{The hire graph of the 3-free door}
\date{}
\author{}

\begin{document}
\maketitle
\thispagestyle{empty}
\vspace{-2.8em}

Let \(\chi_3\) be the non-principal Dirichlet character modulo \(3\):
\(\chi_3(n)=0\) if \(3\mid n\), and \(\chi_3(n)=\bigl(\frac{n}{3}\bigr)\) otherwise,
so \(\chi_3(n)=+1\) if \(n\equiv 1\pmod{3}\) and \(\chi_3(n)=-1\) if \(n\equiv 2\pmod{3}\).
For an odd prime \(p\neq 3\) write
\[
m_0(p)=p+\chi_3(p),
\qquad
m_1(p)=p-\chi_3(p).
\]

\begin{lemma}
\(m_0(p)\) is even and not divisible by \(3\).
\(m_1(p)\) is divisible by \(6\).
\(\{m_0(p),m_1(p)\}=\{p-1,p+1\}\).
In particular \(3\) never divides a \(3\)-free door.
\end{lemma}

\begin{proof}
The residues of \(p-1\), \(p\), \(p+1\) modulo \(3\) are a permutation of \(\{0,1,2\}\), and \(p\not\equiv 0\pmod{3}\).
Hence exactly one neighbour of \(p\) is divisible by \(3\).
Both neighbours are even.
If \(p\equiv 1\pmod{3}\) then \(\chi_3(p)=+1\), so \(m_0(p)=p+1\not\equiv 0\pmod{3}\) and \(m_1(p)=p-1\equiv 0\pmod{3}\).
If \(p\equiv 2\pmod{3}\) then \(\chi_3(p)=-1\), so \(m_0(p)=p-1\not\equiv 0\pmod{3}\) and \(m_1(p)=p+1\equiv 0\pmod{3}\).
In both cases \(m_1(p)\) is even and divisible by \(3\), hence by \(6\).
\end{proof}

\noindent
Examples.
\(p=11\equiv 2\pmod{3}\), so \(\chi_3(11)=-1\), \(m_0(11)=10=2\cdot 5\), \(m_1(11)=12=2^2\cdot 3\).
\(p=13\equiv 1\pmod{3}\), so \(\chi_3(13)=+1\), \(m_0(13)=14=2\cdot 7\), \(m_1(13)=12=2^2\cdot 3\).

\begin{lemma}
The four residue classes of odd primes modulo \(12\) determine the doors as follows.
\end{lemma}

\begin{center}
\small
\begin{tabular}{@{}cccccc@{}}
\toprule
\(p\pmod{12}\) & \(\chi_3\) & \(m_0\pmod{12}\) & \(v_2(m_0)\) & \(m_1\pmod{12}\) & \(v_2(m_1)\) \\
\midrule
\(1\)  & \(+1\) & \(2\)  & \(=1\) & \(0\) & \(\ge 2\) \\
\(5\)  & \(-1\) & \(4\)  & \(\ge 2\) & \(6\) & \(=1\) \\
\(7\)  & \(+1\) & \(8\)  & \(\ge 2\) & \(6\) & \(=1\) \\
\(11\) & \(-1\) & \(10\) & \(=1\) & \(0\) & \(\ge 2\) \\
\bottomrule
\end{tabular}
\end{center}

\begin{proof}
The classes \(\{1,5,7,11\}\pmod{12}\) are the units modulo \(12\).
Each line is the conjunction of Lemma~1 with the unique lifting of that class modulo \(4\) and modulo \(3\).
Residues \(2\) and \(10\) modulo \(12\) are \(2\pmod{4}\), so \(v_2(m_0)=1\); residues \(4\) and \(8\) are \(0\pmod{4}\), so \(v_2(m_0)\ge 2\).
\end{proof}

\begin{lemma}
If \(p\) and \(p+2\) are both prime and \(p>3\), then \(p\equiv 2\pmod{3}\) and
\[
m_1(p)=m_1(p+2)=p+1.
\]
\end{lemma}

\begin{proof}
The pair of residues of \((p,p+2)\) modulo \(3\) is either \((1,0)\) or \((2,1)\).
The first makes \(p+2\) divisible by \(3\), impossible for \(p+2>3\).
Thus \(p\equiv 2\pmod{3}\) and \(p+2\equiv 1\pmod{3}\), so \(\chi_3(p)=-1\) and \(\chi_3(p+2)=+1\).
Then \(m_1(p)=p+1\) and \(m_1(p+2)=p+1\).
\end{proof}

\begin{corollary}
If an even integer is not divisible by \(3\), at most one of its two neighbours is an odd prime greater than \(3\).
\end{corollary}

\begin{lemma}
If \(p\) and \(q=2p+1\) are both prime and \(p>3\), then \(p\equiv 2\pmod{3}\) and \(m_0(q)=2p\).
\end{lemma}

\begin{proof}
If \(p\equiv 1\pmod{3}\) then \(q=2p+1\equiv 0\pmod{3}\), impossible for \(q>3\).
Thus \(p\equiv 2\pmod{3}\), so \(q\equiv 2\pmod{3}\), \(\chi_3(q)=-1\), and \(m_0(q)=q-1=2p\).
\end{proof}

\paragraph{The hire graph.}
An odd prime \(p\neq 3\) is an \emph{owner} of an odd prime \(q\) when \(q\mid m_0(p)\).
By Dirichlet's theorem on the two admissible classes modulo \(3q\),
\[
p\equiv 1\pmod{3},\ \ p\equiv -1\pmod{q}
\qquad\text{and}\qquad
p\equiv 2\pmod{3},\ \ p\equiv 1\pmod{q},
\]
every odd prime \(q\neq 3\) has an owner.
Hence the set of primes that occur as \(2\) or as an odd prime factor of some \(m_0(p)\) is all primes except \(3\).

Spectra are taken on finite windows.
For \(X\ge 5\), let the owners be the odd primes \(p\neq 3\) with \(p\le X\).
Let \(S_X\) consist of \(2\) together with every odd prime factor of \(m_0(p)\) for such owners, and write \(n=|S_X|\).
The \emph{hire graph} \(G_X\) has vertex set \(S_X\), a spoke from \(2\) to every other vertex, and a directed gold arc \(p\to q\) whenever \(p\le X\), \(p,q\in S_X\setminus\{2\}\), and \(q\mid m_0(p)\) (gold witnessed by owners in the window).
The underlying undirected gold graph has an edge \(\{p,q\}\) when at least one of the two arcs is present.
Let \(H_X\) be the combinatorial Laplacian of the undirected graph (spokes plus undirected gold), and let \(L'_X\) be the Laplacian of the undirected gold graph on \(S_X\setminus\{2\}\).

The relation \(q\mid m_0(p)\) is not symmetric:
\(5\mid m_0(11)=10\) while \(11\nmid m_0(5)=4\).

\begin{lemma}
There is a gold arc \(p\to q\) if and only if \(p\equiv -\chi_3(p)\pmod{q}\).
\end{lemma}

\begin{proof}
The congruence is \(q\mid\bigl(p+\chi_3(p)\bigr)\), i.e.\ \(q\mid m_0(p)\).
\end{proof}

\begin{lemma}
Every undirected gold edge \(\{p,q\}\) is the base of a triangle \(2\textrm{--}p\textrm{--}q\) in \(G_X\).
\end{lemma}

\begin{proof}
Both \(p\) and \(q\) are adjacent to \(2\) by construction.
\end{proof}

\begin{lemma}
\(\lambda_{\max}(H_X)=n\).
The vector with value \(n-1\) at \(2\) and value \(-1\) at every other vertex is an eigenvector for this eigenvalue, independently of the gold edges.
\end{lemma}

\begin{proof}
Let \(v\) be that vector.
On the star alone, \(H_Xv=nv\).
Each gold edge has both ends in the set where \(v\) is constant, so contributes \(0\) to \(H_Xv\).
\end{proof}

\begin{lemma}
Let \(0=\mu_1\le\mu_2\le\cdots\le\mu_{n-1}\) be the eigenvalues of \(L'_X\).
Then
\[
\operatorname{spec}(H_X)
=
\{0,\,n\}
\cup
\{1+\mu_2,\ldots,1+\mu_{n-1}\}.
\]
In particular \(\lambda_2(H_X)=1\) if and only if the undirected gold graph on \(S_X\setminus\{2\}\) is disconnected.
\end{lemma}

\begin{proof}
Order the vertices as \(2\) followed by \(S_X\setminus\{2\}\).
Then
\[
H_X
=
\begin{pmatrix}
k & -\mathbf{1}^{\mathsf T}\\
-\mathbf{1} & I+L'_X
\end{pmatrix},
\]
where \(k=n-1\).
Vectors constant on \(S_X\setminus\{2\}\) yield eigenvalues \(0\) and \(n\).
On the complement \(\mathbf{1}^{\mathsf T}y=0\) the lower block acts as \(1+\mu\), and the copy of \(\ker L'_X\) spanned by \(\mathbf{1}\) has already been used.
The remaining eigenvalues of \(L'_X\) are \(\mu_2,\ldots,\mu_{n-1}\).
\end{proof}

\noindent
On the windows below, undirected gold remains disconnected, so \(\lambda_2(H_X)=1\).

\begin{center}
\small
\begin{tabular}{@{}rrrrccc@{}}
\toprule
\(X\) & \(|S_X|\) & gold edges & gold arcs & \(L'_X\) connected & \(\mu_2(L'_X)\) & \(\lambda_2(H_X)\) \\
\midrule
\(30\)  & \(4\)  & \(1\)  & \(1\)  & no & \(0\) & \(1\) \\
\(50\)  & \(6\)  & \(3\)  & \(3\)  & no & \(0\) & \(1\) \\
\(100\) & \(12\) & \(7\)  & \(7\)  & no & \(0\) & \(1\) \\
\(200\) & \(18\) & \(13\) & \(13\) & no & \(0\) & \(1\) \\
\bottomrule
\end{tabular}
\end{center}

\noindent
\(X=30\): \(S=\{2,5,7,11\}\), gold \(\{5,11\}\).
\(X=50\): \(S=\{2,5,7,11,19,23\}\), gold \(\{5,11\}\), \(\{5,19\}\), \(\{11,23\}\).
\(X=100\): gold \(\{5,11\}\), \(\{5,19\}\), \(\{5,41\}\), \(\{7,13\}\), \(\{7,29\}\), \(\{11,23\}\), \(\{19,37\}\).

\paragraph{Variant: the sluice.}
The hire graph \(G_X\) uses \(m_0\) throughout.
A finite set \(\Sigma\) of odd primes defines a variant door:
the \(\Sigma\)-door of \(p\) is \(m_0(p)\) if no prime of \(\Sigma\) divides \(m_0(p)\); otherwise \(m_1(p)\) if no prime of \(\Sigma\) divides \(m_1(p)\); and is otherwise undefined.

\begin{lemma}
If \(\Sigma=\{5\}\), the \(\Sigma\)-door is defined for every odd prime \(p\neq 3\).
It differs from \(m_0(p)\) if and only if \(p\equiv 11\) or \(19\pmod{30}\).
\end{lemma}

\begin{proof}
The condition \(5\mid m_0(p)\) is \(p\equiv -\chi_3(p)\pmod{5}\).
Combined with \(p\) odd and the value of \(\chi_3(p)\), this is the pair of classes \(11,19\pmod{30}\).
On those classes \(m_1(p)\) is divisible by \(6\) and not by \(5\): the two neighbours of \(p\) differ by \(2\), so at most one is divisible by \(5\).
\end{proof}

\begin{lemma}
If \(\Sigma=\{5,7\}\), the \(\Sigma\)-door is undefined if and only if
\[
p\equiv 29,\ 41,\ 169,\ \text{or }181\pmod{210}.
\]
These are four residue classes modulo \(210\), of density \(4/\varphi(210)=1/12\) among units modulo \(210\).
\end{lemma}

\begin{proof}
The door is undefined precisely when \(5\) divides one neighbour of \(p\) and \(7\) divides the other, i.e.\
\[
\bigl\{p-1,\,p+1\bigr\}
\equiv
\bigl\{0\pmod{5},\,0\pmod{7}\bigr\}
\]
as sets.
The system \(p\equiv \pm 1\pmod{5}\), \(p\equiv \mp 1\pmod{7}\), together with \(p\) odd, cuts four classes modulo \(210=2\cdot 3\cdot 5\cdot 7\).
The odd representatives are \(29\), \(41\), \(169\equiv -41\), and \(181\equiv -29\).
There are \(\varphi(210)=48\) units modulo \(210\), so the density among those units is \(4/48=1/12\).
\end{proof}

\end{document}
